\documentclass[11pt,letterpaper]{article}

\usepackage[margin=1in]{geometry}
\usepackage[T1]{fontenc}
\usepackage{lmodern}
\usepackage{microtype}
\usepackage{enumitem}

\setcounter{page}{6}
\setlength{\parindent}{0pt}
\setlength{\parskip}{5pt}
\setlist{nosep,leftmargin=1.6em}

\begin{document}

\begin{center}
  {\Large\bfseries CRC Package Interface Specification}\\[3pt]
  {\large Model-matrix implementation addendum}
\end{center}

\section{Model-matrix infrastructure}

The package now contains initial infrastructure for converting parsed model
specifications into complete cell grids and observation-level latent states.
The implementation remains separate from \texttt{crc\_fit()} until the
remaining matrices and their alignment rules are defined.

\subsection{Capture-model matrix}

\texttt{build\_capture\_matrix()} constructs the Cartesian product of all
capture histories, categorical levels used by the capture model, and latent
classes. For misclassified variables, true-category levels and their ordering
are obtained from the corresponding confusion matrices. The function returns:

\begin{itemize}
  \item \texttt{cells}, the complete capture-model cell grid;
  \item \texttt{matrix}, the design matrix generated from the parsed terms;
  \item \texttt{unobserved}, an indicator of all-zero capture histories.
\end{itemize}

With $J$ capture sources, $K$ latent classes, and categorical variables with
$C_1,\ldots,C_p$ levels, the grid contains
$2^J K\prod_{q=1}^p C_q$ rows.

\subsection{Initial observation states}

\texttt{build\_observation\_states()} repeats every observed unit once for
each latent class. It retains the variables needed by the capture model,
outcome model, and misclassification mechanisms. It also keeps aligned copies
of the standardized outcome and observed values of misclassified variables.
An observation identifier maps every expanded row back to its original unit.

\subsection{Misclassified-variable expansion}

\texttt{expand\_misclass\_states()} expands the current states over every
possible true level of one misclassified variable. Each expanded row receives
the contribution
\[
  \Pr(\widetilde G_j=h\mid G_j=g)
\]
from its confusion matrix. When \texttt{true\_if} applies, the observed
category receives probability one and all other candidate categories receive
probability zero.

For several misclassified variables, separate confusion matrices are used and
their contributions are multiplied. This corresponds to conditional
independence of the misclassification mechanisms given their respective true
categories. Dependent mechanisms may be supported in a future extension.
These contributions are likelihood terms, not posterior probabilities, and
must not be normalized at this stage.

\newpage

\section{Required next steps}

The model-matrix work should continue in the following order.

\begin{enumerate}
  \item Add a small wrapper that calls
  \texttt{expand\_misclass\_states()} for every variable in
  \texttt{misclass}. When no mechanism is specified, it should assign a
  contribution of one to every state.

  \item Build the outcome-model design matrix from the fully expanded state
  data. It must remain row-aligned with the standardized outcome,
  observation identifiers, and misclassification contributions.

  \item Define indices linking each observation state to the corresponding
  row of the complete capture-cell matrix. This avoids repeatedly matching
  capture histories, true categorical values, and latent classes during the
  EM iterations.

  \item Create a single internal model-matrix object containing the capture
  grid, capture design matrix, expanded observation states, outcome design
  matrix, alignment indices, and outcome metadata.

  \item Integrate construction of this object into \texttt{crc\_fit()} while
  retaining the current unfitted-object behavior.

  \item Add focused tests for dimensions, row alignment, factor-level order,
  all-zero capture cells, multiple misclassified variables, and
  \texttt{true\_if} restrictions.
\end{enumerate}

Only after these structures and invariants are established should development
proceed to parameter initialization and the EM algorithm. In the E-step, the
misclassification contributions will be combined with capture-model expected
cell frequencies and outcome likelihood contributions, and then normalized
within each observed unit to obtain posterior state probabilities.

\end{document}
